% sec3_4.tex — Section 3.4: Generalized Method of Moments Defined

The orthogonality conditions state that a set of population moments are all equal
to zero.  The basic principle of the \textbf{method of moments} is to choose the parameter
estimate so that the corresponding sample moments are also equal to zero.
The population moments in the orthogonality conditions are
$\E[\mathbf{g}(\mathbf{w}_i;\, \boldsymbol{\delta})]$.  Its sample analogue is the sample mean of
$\mathbf{g}(\mathbf{w}_i;\, \boldsymbol{\delta})$ (where $\mathbf{g}(\mathbf{w}_i;\, \boldsymbol{\delta})$ is
defined in \eqref{eq:3.3.1}) evaluated at some hypothetical value $\tilde{\boldsymbol{\delta}}$ of
$\boldsymbol{\delta}$:
\begin{equation}
  \underset{(K \times 1)}{\mathbf{g}_n(\tilde{\boldsymbol{\delta}})}
  \equiv \frac{1}{n} \sum_{i=1}^{n} \mathbf{g}(\mathbf{w}_i;\, \tilde{\boldsymbol{\delta}}).
  \label{eq:3.4.1}
\end{equation}

Applying the method of moments principle to our model amounts to choosing
a $\tilde{\boldsymbol{\delta}}$ that solves the system of $K$ simultaneous equations in $L$ unknowns:
$\mathbf{g}_n(\tilde{\boldsymbol{\delta}}) = \mathbf{0}$, which is the sample analog of \eqref{eq:3.3.3}.
Because the estimation equation is linear, $\mathbf{g}_n(\tilde{\boldsymbol{\delta}})$ can be written as
\begin{align}
  \mathbf{g}_n(\tilde{\boldsymbol{\delta}})
  &= \frac{1}{n} \sum_{i=1}^{n} \mathbf{x}_i \cdot
     (y_i - \mathbf{z}_i' \tilde{\boldsymbol{\delta}})
  \quad \text{(by definition \eqref{eq:3.3.1} of
    $\mathbf{g}(\mathbf{w}_i;\, \boldsymbol{\delta})$)} \notag \\
  &= \frac{1}{n} \sum \mathbf{x}_i \cdot y_i
     - \Bigl(\frac{1}{n} \sum \mathbf{x}_i \mathbf{z}_i'\Bigr)
       \tilde{\boldsymbol{\delta}}
  \equiv \mathbf{s}_{\mathbf{xy}} - \mathbf{S}_{\mathbf{xz}}\,
         \tilde{\boldsymbol{\delta}},
  \label{eq:3.4.2}
\end{align}
where $\mathbf{s}_{\mathbf{xy}}$ and $\mathbf{S}_{\mathbf{xz}}$ are the corresponding sample moments
of $\boldsymbol{\sigma}_{\mathbf{xy}}$ and $\boldsymbol{\Sigma}_{\mathbf{xz}}$:
\[
  \underset{(K \times 1)}{\mathbf{s}_{\mathbf{xy}}}
  \equiv \frac{1}{n} \sum_{i=1}^{n} \mathbf{x}_i \cdot y_i
  \quad \text{and} \quad
  \underset{(K \times L)}{\mathbf{S}_{\mathbf{xz}}}
  \equiv \frac{1}{n} \sum_{i=1}^{n} \mathbf{x}_i \mathbf{z}_i'.
\]
So the sample analog $\mathbf{g}_n(\tilde{\boldsymbol{\delta}}) = \mathbf{0}$ is a system of $K$
linear equations in $L$ unknowns:
\begin{equation}
  \mathbf{S}_{\mathbf{xz}}\, \tilde{\boldsymbol{\delta}}
  = \mathbf{s}_{\mathbf{xy}}.
  \label{eq:3.4.3}
\end{equation}
This is the sample analogue of \eqref{eq:3.3.4}.  If there are more orthogonality conditions
than parameters, that is, if $K > L$, then the system may not have a solution.  The
extension of the method of moments to cover this case as well is the \textbf{generalized
method of moments (GMM)}.

%% ─── Method of Moments ─────────────────────────────────────────────────────
\subsection*{Method of Moments}

If the equation is exactly identified, then $K = L$ and $\boldsymbol{\Sigma}_{\mathbf{xz}}$ is square and
invertible.  Since under Assumption~3.2 $\mathbf{S}_{\mathbf{xz}}$ converges to
$\boldsymbol{\Sigma}_{\mathbf{xz}}$ almost surely, $\mathbf{S}_{\mathbf{xz}}$ is invertible for sufficiently
large sample size $n$ with probability one.  Thus, when the sample size is large, the
system of simultaneous equations \eqref{eq:3.4.3} has a unique solution given by
\begin{equation}
  \hat{\boldsymbol{\delta}}_{\text{IV}}
  = \mathbf{S}_{\mathbf{xz}}^{-1}\, \mathbf{s}_{\mathbf{xy}}
  = \Bigl(\frac{1}{n} \sum_{i=1}^{n} \mathbf{x}_i \mathbf{z}_i'\Bigr)^{-1}
    \frac{1}{n} \sum_{i=1}^{n} \mathbf{x}_i \cdot y_i.
  \label{eq:3.4.4}
\end{equation}
This estimator is also called the \textbf{instrumental variables estimator} with $\mathbf{x}_i$
serving as instruments.  Because it is defined for the exactly identified case, the formula
assumes that there are as many instruments as regressors.  If, moreover,
$\mathbf{z}_i = \mathbf{x}_i$ (all the regressors are predetermined or orthogonal to the error term),
then $\hat{\boldsymbol{\delta}}_{\text{IV}}$ reduces to the OLS estimator.  Thus, the OLS estimator is a method of
moments estimator.

%% ─── Generalized Method of Moments ─────────────────────────────────────────
\subsection*{Generalized Method of Moments}

If the equation is overidentified so that $K > L$, we cannot in general choose an
$L$-dimensional $\tilde{\boldsymbol{\delta}}$ to satisfy the $K$ equations in \eqref{eq:3.4.3}.
If we cannot set $\mathbf{g}_n(\tilde{\boldsymbol{\delta}})$ exactly equal to $\mathbf{0}$, we can at
least choose $\tilde{\boldsymbol{\delta}}$ so that
$\mathbf{g}_n(\tilde{\boldsymbol{\delta}})$ is as close to $\mathbf{0}$ as possible.  To
make precise what we mean by ``close,'' we define the distance between any two
$K$-dimensional vectors $\boldsymbol{\xi}$ and $\boldsymbol{\eta}$ by the quadratic form
$(\boldsymbol{\xi} - \boldsymbol{\eta})' \hat{\mathbf{W}}
(\boldsymbol{\xi} - \boldsymbol{\eta})$, where $\hat{\mathbf{W}}$, sometimes called the
\textbf{weighting matrix}, is a symmetric and positive definite
matrix defining the distance.\footnote{Do not let the word ``weight'' confuse you between
GMM and weighted least squares (WLS).  In GMM, the weighting is applied to the sample
mean $\bar{\mathbf{g}}$, while in WLS it applies to each observation.}

\begin{definition}[GMM estimator]
\label{def:3.1}
\emph{Let $\hat{\mathbf{W}}$ be a $K \times K$ symmetric positive definite
matrix, possibly dependent on the sample, such that
$\hat{\mathbf{W}} \pto \mathbf{W}$ as the sample size $n$ goes to infinity with $\mathbf{W}$
symmetric and positive definite.  The \textbf{GMM estimator} of $\boldsymbol{\delta}$, denoted
$\hat{\boldsymbol{\delta}}(\hat{\mathbf{W}})$, is}
\begin{equation}
  \hat{\boldsymbol{\delta}}(\hat{\mathbf{W}})
  \equiv \argmin_{\tilde{\boldsymbol{\delta}}}\;
  J(\tilde{\boldsymbol{\delta}},\, \hat{\mathbf{W}}),
  \label{eq:3.4.5}
\end{equation}
\emph{where}
\[
  J(\tilde{\boldsymbol{\delta}},\, \hat{\mathbf{W}})
  \equiv n \cdot
  \mathbf{g}_n(\tilde{\boldsymbol{\delta}})'\,
  \hat{\mathbf{W}}\,
  \mathbf{g}_n(\tilde{\boldsymbol{\delta}}).
\]
\end{definition}

(The reason the distance
$\mathbf{g}_n(\tilde{\boldsymbol{\delta}})'\hat{\mathbf{W}}\,\mathbf{g}_n(\tilde{\boldsymbol{\delta}})$
is multiplied by the sample size ($n$) becomes clear in Section~3.6.)
The weighting matrix is allowed to be random and depend on the sample size, to cover
the possibility that the matrix is estimated from the sample.  The definition makes clear
that the GMM is a special case of \textbf{minimum distance estimation}.  In minimum distance
estimation, $\plim\, \mathbf{g}_n(\boldsymbol{\delta}) = \mathbf{0}$, as here, but the
$\mathbf{g}_n(\cdot)$ function is not necessarily a sample mean.

Since $\mathbf{g}_n(\tilde{\boldsymbol{\delta}})$ is linear in $\tilde{\boldsymbol{\delta}}$ (see
\eqref{eq:3.4.2}), the objective function is quadratic in $\tilde{\boldsymbol{\delta}}$ when the
equation is linear:
\begin{equation}
  J(\tilde{\boldsymbol{\delta}},\, \hat{\mathbf{W}})
  = n \cdot (\mathbf{s}_{\mathbf{xy}} - \mathbf{S}_{\mathbf{xz}}\,
    \tilde{\boldsymbol{\delta}})'\,
    \hat{\mathbf{W}}\,
    (\mathbf{s}_{\mathbf{xy}} - \mathbf{S}_{\mathbf{xz}}\,
    \tilde{\boldsymbol{\delta}}).
  \label{eq:3.4.6}
\end{equation}
It is left to you to show that the first-order condition for the minimization of this
with respect to $\tilde{\boldsymbol{\delta}}$ is
\begin{equation}
  \underset{(L \times K)}{\mathbf{S}_{\mathbf{xz}}'} \;
  \underset{(K \times K)}{\hat{\mathbf{W}}} \;
  \underset{(K \times 1)}{\mathbf{s}_{\mathbf{xy}}}
  = \underset{(L \times K)}{\mathbf{S}_{\mathbf{xz}}'} \;
    \underset{(K \times K)}{\hat{\mathbf{W}}} \;
    \underset{(K \times L)}{\mathbf{S}_{\mathbf{xz}}} \;
    \underset{(L \times 1)}{\tilde{\boldsymbol{\delta}}}\,.
  \label{eq:3.4.7}
\end{equation}
If Assumptions~3.2 and 3.4 hold, then $\mathbf{S}_{\mathbf{xz}}$ is of full column rank for sufficiently
large $n$ with probability one.  Then, since $\hat{\mathbf{W}}$ is positive definite, the
$L \times L$ matrix $\mathbf{S}_{\mathbf{xz}}' \hat{\mathbf{W}} \mathbf{S}_{\mathbf{xz}}$ is nonsingular.
So the unique solution can be obtained by multiplying both sides by the inverse of
$\mathbf{S}_{\mathbf{xz}}' \hat{\mathbf{W}} \mathbf{S}_{\mathbf{xz}}$.  That unique solution is the
GMM estimator:
\begin{equation}
\boxed{%
  \text{GMM estimator:} \quad
  \hat{\boldsymbol{\delta}}(\hat{\mathbf{W}})
  = (\mathbf{S}_{\mathbf{xz}}' \hat{\mathbf{W}} \mathbf{S}_{\mathbf{xz}})^{-1}\,
    \mathbf{S}_{\mathbf{xz}}' \hat{\mathbf{W}}\, \mathbf{s}_{\mathbf{xy}}.
}
  \label{eq:3.4.8}
\end{equation}
If $K = L$, then $\mathbf{S}_{\mathbf{xz}}$ is a square matrix and \eqref{eq:3.4.8} reduces to the IV
estimator \eqref{eq:3.4.4}.  The GMM is indeed a generalization of the method of moments.

%% ─── Sampling Error ────────────────────────────────────────────────────────
\subsection*{Sampling Error}

For later use, we derive the expression for the sampling error.  Multiplying both
sides of the estimation equation $y_i = \mathbf{z}_i' \boldsymbol{\delta} + \varepsilon_i$ from left
by $\mathbf{x}_i$ and taking the averages yields
\begin{equation}
  \mathbf{s}_{\mathbf{xy}} = \mathbf{S}_{\mathbf{xz}}\, \boldsymbol{\delta}
  + \bar{\mathbf{g}},
  \label{eq:3.4.9}
\end{equation}
where
\begin{equation}
  \bar{\mathbf{g}} \equiv \frac{1}{n} \sum_{i=1}^{n} \mathbf{x}_i \cdot \varepsilon_i
  = \frac{1}{n} \sum_{i=1}^{n} \mathbf{g}(\mathbf{w}_i;\, \boldsymbol{\delta})
  = \mathbf{g}_n(\boldsymbol{\delta}).
  \label{eq:3.4.10}
\end{equation}
Substituting \eqref{eq:3.4.9} into \eqref{eq:3.4.8}, we obtain
\begin{equation}
  \hat{\boldsymbol{\delta}}(\hat{\mathbf{W}}) - \boldsymbol{\delta}
  = (\mathbf{S}_{\mathbf{xz}}' \hat{\mathbf{W}} \mathbf{S}_{\mathbf{xz}})^{-1}\,
    \mathbf{S}_{\mathbf{xz}}' \hat{\mathbf{W}}\, \bar{\mathbf{g}}.
  \label{eq:3.4.11}
\end{equation}

%% ─── Questions for Review ──────────────────────────────────────────────────
\bigskip
\begin{center}\rule{0.6\textwidth}{0.4pt}\end{center}
\noindent\textsc{Questions for Review}
\medskip

\begin{enumerate}
\item Verify that \eqref{eq:3.1.12} is the IV estimator of $\alpha_1$ when the method of moments is
  applied to the demand equation \eqref{eq:3.1.2a} with $(1, x_i)$ as instruments.

\item If the equation is just identified, what is the minimized value of
  $J(\tilde{\boldsymbol{\delta}},\, \hat{\mathbf{W}})$?

\item Even if the equation is overidentified, the population version, \eqref{eq:3.3.4}, of the
  system of $K$ equations \eqref{eq:3.4.3} has a solution.  Then, why does not the sample
  version, \eqref{eq:3.4.3}, have a solution?  \textbf{Hint:} The population version has a solution
\end{enumerate}
